%%%%%%%%%%%%%%%%%%%%%%%%%%%%%%%%%%%%%%%%%
% Template for ASE NTU Qualify Exam
% Version 1.0 (26-Jan-2022)
%
% Original author:
% Adriana Lopes dos Santos and Daniel Vaulot
%%%%%%%%%%%%%%%%%%%%%%%%%%%%%%%%%%%%%%%%%

\documentclass[a4paper,10pt]{article}
\usepackage[left=2cm,top=1.27cm,right=2cm,bottom=1.27cm]{geometry}
\usepackage{dmasproject} %use for bib style and all style set up
% if you need additional LaTeX packages, add them here
\graphicspath{{Fig/}}
\usepackage{setspace}
\usepackage{soul} 
\usepackage{gensymb}
\usepackage{hyperref}
\hypersetup{colorlinks = true, citecolor = blue, linkcolor = blue, urlcolor = blue} % remove color for links 
\usepackage{lipsum}
%\usepackage{titlesec}			% to reduce the space before and after the paragraph and the left spacing
\usepackage{fancyhdr}
\usepackage{multicol}



\usepackage{siunitx}
    \usepackage[compact]{titlesec}
    \titlespacing{\section}{0pt}{0.5ex}{0ex}
    \titlespacing{\subsection}{0pt}{0.5ex}{0ex}
    \titlespacing{\subsubsection}{0pt}{0.5ex}{0ex}
\renewcommand{\baselinestretch}{1.0}
\setlength{\parindent}{0pt}
\setlength{\parskip}{0.2cm}
\titleformat{\section}
  {\centering\normalfont\fontsize{12}{14}\bfseries}{\thesection}{1em}{}
 \titleformat{\subsection}
  {\normalfont\fontsize{10}{15}\bfseries}{\thesubsection}{1em}{}
\newcommand{\squeezeup}{\vspace{-2.5mm}}
\usepackage{mathptmx}
\pagestyle{fancy}% sets both header and footer to nothing
\renewcommand{\headrulewidth}{0pt}
\fancyhf{}

\makeatletter
\renewcommand{\paragraph}{%
  \@startsection{paragraph}{4}%
  {\z@}{0ex \@plus 1ex \@minus .2ex}{-1em}%
  {\normalfont\normalsize\bfseries}%
}
\makeatother

\pagestyle{fancy}
\fancyhf{}
\rhead{Title of Your Paper - p. \thepage}
\lhead{}
\cfoot{}

%\fancyfoot[L]{DOI: 10.2112/SI65-xxx.1 received 05 April 2020;\newline accepted Day Month 202X}


\begin{document}


 	\noindent\begin{minipage}[h]{1\textwidth}
 	
 	\Large
    \textbf{Title of your paper}
        
    \vspace{0.8cm}
 	\end{minipage}
 	\begin{minipage}[b]{0.5\textwidth}
    \textbf{Surname, First Name\textsuperscript{1}}\par
        \vspace{0.4cm}
    \small
    \textsuperscript{1} Asian School of the Environment, Nanyang Technological\\University, Singapore simw0034@ntu.edu.sg\par

 \end{minipage} 
   \hfill
\begin{minipage}[b]{0.4\textwidth}
 	\noindent\Large
    \centering\textbf{ASIAN SCHOOL OF\\THE ENVIRONMENT}
\end{minipage} 

\vspace{1.2cm}

\hspace{1cm}
\hfill
\begin{minipage}[h]{0.92\textwidth}
\textbf{ABSTRACT}

An abstract is NOT another introduction. Its purpose is to enable potential reader determine whether your work is relevant to the work he or she is undertaking. When searching the Index and Abstract data bases of the library to acquire suitable material to start your literature research you will have used abstracts for the same purpose. Most likely you will have wanted the author to keep it concise and to the point. Abstract may be written in structured or unstructured forms. In their most basic form, structured abstracts are based on following headings where possible: Aims and objectives; Background, stating what is already known about the topic. The abstract should consist of a condensed summary of the findings of the research. The abstract should accurately reflect the content of the paper. The abstract should not include references or abbreviations. The abstract should be followed by two to six keywords, which accurately identify the paper’s subject, purpose and focus. Abstracts should emphasize new and important aspects of the study or observations. Authors must use this document as a basis of their articles. Abstract should not exceed 300 words 
 
     \vspace{0.4cm}
\noindent
Surname, F.N., 2022. Title of Your paper. \textit{Proceedings of the 2021 Qualifying Exam Series, Asian School of the Environment Special Publication.}

    \vspace{0.4cm}
\noindent\textbf{ADDITIONAL INDEX WORDS}:  Interesting key words or phrases, that will, get your, paper cited; 

\rule{1\textwidth}{0.5pt}
\end{minipage}

\thispagestyle{empty}

% 5 page including figures. Abstract and references not included. No supplementary figures.
\vspace{7.6cm}
DOI: 10.2112/SI65-xxx.1 received 01 November 2021;\newline accepted Day Month 2021

\newpage
\begin{multicols}{2}
\section*{INTRODUCTION}
The introduction basically explains to the reader in more detail some or all of the following points in a sentence or paragraph each:
Why you have undertaken the research and written this paper?
What has been done before (this commonly takes a few paragraphs). This commonly requires a suscinct informative statement that demonstrates that you are fully aware of the context of your study and can demonstrate that you have adequately considered all the most relevant literature (not just your supervisor’s papers) on the topic \parencite{Author1}.
What is the scope of the investigation, i.e. what is covered and what is not. Are there any limitations inherent in the study. Your QE is a small short study and it is commonly a part of a largere program. Explain how your work fits in the bigger picture then zoom in on what you did and why its important
The theories that form the context for this study (studies with a heavy literature review component may need a separate literature review sub-section)
An outline of the structure of your paper. 

\section*{METHODS}
\paragraph*{Study Location.}
In the second section of your paper focus in on what you did (Table~\ref{table:table_1}). If you did fieldwork describe the setting with a study site description before a second section with. (67.4797 \degree N, 63.7895 \degree, Figure~\ref{fig:figure1}). 


\section*{RESULTS}

The function of this section is to objectively present your key results, without interpretation, in an orderly and logical sequence. Commonly this is done using both text and illustrative materials (Tables and Figures). The result section always begins with some text, reporting the key results and referring to your figures and tables as you proceed. Summaries of any statistical analyses may appear either in the text (usually parenthetically) or in the relevant Tables or Figures (in the legend or as footnotes to the Table or Figure). The Results section should be organized around Tables and/or Figures. Any tables and figures should be sequenced to present your key findings in a logical order. The text of the Results section should be crafted to follow the sequence of tables and figures and highlight the evidence needed to answer the questions/hypotheses you investigated. Important negative results should be reported, too. Authors usually write the text of the results section based upon the sequence of Tables and Figures (Figure~\ref{fig:figure2}).

In terms of style try to write the text of the Results section concisely and objectively. The passive voice will likely dominate here, but use the active voice as much as possible. Use the past tense. Avoid repetitive paragraph structures. Do not interpret the data here. The transition into interpretive language can be a slippery slope (Figure~\ref{fig:figure3}).

\section*{DISCUSSION}
The Discussion is where you interpret your results in light of what was already known about the topic, and to explain our new understanding of the problem after taking your results into consideration. A good Discussion should always connect to the Introduction by way of the question(s) or hypotheses you posed and the literature you cited, but it does not simply repeat or rearrange the Introduction. Instead, it tells how your study has moved us forward from the place you left us at the end of the Introduction.

\paragraph*{Fundamental questions to answer here include:}
\begin{itemize}
    \item Do your results provide answers to your testable hypotheses? If so, how do you interpret your findings?
    \item Do your findings agree with what others have shown? If not, do they suggest an alternative explanation or perhaps a unforseen design flaw in your experiment (or theirs?)
    \item Given your conclusions, what is our new understanding of the problem you investigated and outlined in the Introduction?
    \item If warranted, what would be the next step in your study, e.g., what experiments would you do next?
\end{itemize}

In your discussion it is important to use active voice whenever possible in this section. Watch out for wordy phrases; be concise and make your points clearly. Use of the first person is okay, but too much use of the first person may actually distract the reader from the main points.
Try to organize the Discussion to address each of the experiments or studies for which you presented results; discuss each in the same sequence as presented in the Results, providing your interpretation of what they mean in the larger context of the problem. Do not waste entire sentences restating your results; if you need to remind the reader of the result to be discussed, use "bridge sentences" that relate the result to the interpretation:
"The slow response of the pioneering vegetation suggests that...[interpretation]".
You will necessarily make reference to the findings of others in order to support your interpretations. Use subheadings, if need be, to help organize your presentation. Be wary of mistaking the reiteration of a result for an interpretation, and make sure that no new results are presented here that rightly belong in the results.
You must relate your work to the findings of other studies - including previous studies you may have done and those of other investigators. As stated previously, you may find crucial information in someone else's study that helps you interpret your own data, or perhaps you will be able to reinterpret others' findings in light of yours. In either case you should discuss reasons for similarities and differences between yours and others' findings. Consider how the results of other studies may be combined with yours to derive a new or perhaps better substantiated understanding of the problem. Be sure to state the conclusions that can be drawn from your results in light of these considerations. You may also choose to briefly mention further studies you would do to clarify your working hypotheses. Make sure to reference any outside sources as shown in the Introduction section.
Do not introduce new results in the Discussion. Although you might occasionally include in this section tables and figures which help explain something you are discussing, they must not contain new data (from your study) that should have been presented earlier. They might be flow diagrams, accumulation of data from the literature, or something that shows how one type of data leads to or correlates with another, etc.

\section*{CONCLUSIONS}
Conclusions are often the most difficult part to write, and many writers feel they have nothing left to say after having written the paper. However, you need to keep in mind that most readers read the abstract and conclusion first. A conclusion is where you summarize the paper’s findings and generalize their importance, discuss ambiguous data, and recommend further research. An effective conclusion should provide closure for a paper, leaving the reader feeling satisfied that the concepts have been fully explained. 

\begin{table}[H]
\centering
\caption{All table captions should go above the table. Small tables should be half page. Tables should be centred and only have minimal borders so as to discriminate data from descriptors}
\label{table:table_1}
\resizebox{0.48\textwidth}{!}{
\begin{tabular}{rlcrrr}
  \hline
 \textbf{A} & \textbf{B} & \textbf{C} \\ 
  \hline
\textbf{Item 1} & 0.23 ± 0.19 & 0.61 ± 0.22 & 0.23 ± 0.23 &  \\ 
\textbf{Item 2} & 0.24 ± 0.27 & 2.84 ± 1.98 & 1.36 ± 1.02 &  \\ 
\textbf{Item 3} & 40.84 ± 4.72 & 64.53 ± 37.49 & 27.60 ± 9.24  \\ 
   \hline
\end{tabular}
}
\end{table}



\end{multicols}



\newpage
\section*{FIGURES}


\begin{figure}[H]
    \centering
 	\begin{minipage}[b]{0.3\textwidth}
     	\includegraphics[width=1\linewidth]{figs/study_site.jpg}
     	\caption{Study location \\in XXX.}
     	\label{fig:figure1}
   \end{minipage} 
    \begin{minipage}[h]{0.05\textwidth}
  \end{minipage} 
    \begin{minipage}[b]{0.65\textwidth}
     	\includegraphics[width=1\linewidth]{figs/sampling.jpg}
     	\caption{Sampling Site}
     	\label{fig:figure2}
   \end{minipage}  
  \hfill
\end{figure}
\squeezeup

\begin{figure}[H]
\centering
\includegraphics[width=\textwidth]{figs/temp_nitrate.png}
\vspace{-2em}
\caption{All figure captions should go below the figure. Small figures should be half page.}
\vspace{-1em}
\label{fig:figure3}
\end{figure}





% This will print you references, please do not change it.
\newpage

\printbibliography[title={Literature Cited}]

\end{document}
